\section{Stress tests et architecture v2 --- évaluation indépendante}
\label{sec:robustness}

L'AUROC $= 0.973 \pm 0.012$ (5 graines) ou $= 0.987 \pm 0.011$ (10 graines, config optimisée)
reporté pour H3 (section~\ref{sec:h3}) mesure la prédiction des \emph{labels cluster produits
par EWAT lui-même} depuis l'embedding STGCN. Cette évaluation est \textbf{circulaire} : la cible
($C_i$) est dérivée du même pipeline (encodeur + siamois) qu'on évalue. Le maître de stage a
soulevé cette critique~; nous l'avons investiguée à travers deux phases d'expériences
(stress tests Phase~A et architecture v2 Phase~B/C) qui transforment la critique en
transparence méthodologique et fournissent un \textbf{headline défensif honnête}.

\subsection{Phase A --- 5 stress tests de robustesse}
\label{sec:stress_tests}

Scripts disponibles dans \texttt{experiments/h3\_robustness/}. Synthèse complète dans
\texttt{experiments/h3\_robustness/results.md}.

\begin{table}[h]
\centering
\caption{Verdict synthétique des 5 stress tests sur ewat\_v3.}
\label{tab:stress_tests}
\begin{tabularx}{\linewidth}{l X l}
\toprule
\textbf{Test} & \textbf{Question} & \textbf{Verdict} \\
\midrule
\textbf{A1} distant-window & Le précurseur fait-il de la précursion temporelle sur labels EWAT ? &
  \xmark{} $\Delta($far\,$-$\,near$)=-0.007$ (fuite signature scénario) \\
\textbf{A2} LOSO precursor-only & Généralisation à un type inédit ? &
  \xmark{} top-1 held-out $=0.51 \pm 0.38$ (polarisé) \\
\textbf{A3} permutation null & Signal réel ou bruit ? &
  \checkmark{} AUROC $= 0.893$ vs null $0.49 \pm 0.10$, $p < 0.01$ \\
\textbf{A4} filtre $n_\text{pos} \geq 5$ & Clusters statistiquement reportables ? &
  5/10 (AUROC moyen reportable $= 0.975 \pm 0.020$) \\
\textbf{A5} paired $\Delta(B_4 - B_3)$ & STGCN apporte-t-il du gain prédictif ? &
  \xmark{} $\Delta = +0.005$, IC 95\,\% $= [-0.031, +0.044]$ (contient 0) \\
\bottomrule
\end{tabularx}
\end{table}

\textbf{Verdict synthétique Phase~A.} EWAT discrimine correctement les scénarios actifs
(cohérent avec A3, $p < 0.01$ et H1 silhouette $= 0.78$) mais ne fait pas de précursion
temporelle sur labels EWAT (A1) et ne généralise pas à un type inédit (A2). L'encodeur STGCN
n'apporte pas de gain prédictif agrégé vs un LR sur features brutes (A5). Le headline $0.987$
mesure la cohérence interne du clustering, pas la précursion.

\subsection{Phase B --- Architecture v2 : instance norm + cible Chaos Mesh}
\label{sec:archv2}

Scripts disponibles dans \texttt{experiments/architecture\_v2/}. Synthèse dans
\texttt{experiments/architecture\_v2/results.md}.

\paragraph{Pivot méthodologique.} Remplacer la cible (cluster EWAT auto-référent) par les
\textbf{15 scénarios Chaos Mesh} (vérité terrain indépendante du pipeline). De plus, normaliser
chaque épisode par sa propre moyenne/std sur le régime normal (\emph{instance normalization})
pour éliminer les baselines absolus des services et exposer la dynamique pré-injection.

\paragraph{Diagnostic clé (B1).} Sur ewat\_v4\_strat (épisodes longs 47--51 steps) :
$\Delta($far\,$-$\,near$)= -0.063$ avec instance norm, $-0.043$ avec global norm sur cible
Chaos Mesh. Donc \textbf{il existe bien une dynamique pré-injection} dans le signal --- la
fuite signature scénario observée en A1 ($\Delta \approx 0$) était un artefact de la
circularité des labels EWAT, pas du signal.

\paragraph{Correction critique du split ewat\_v4.}
Le split temporel original de ewat\_v4 avait 4 scénarios entièrement absents du training
(\texttt{faulty\_deploy\_overlap}, \texttt{noisy\_neighbor}, \texttt{memory\_pressure},
\texttt{resource\_leak}). Cela conduisait à un test AUROC trivial de $0.500$. Un nouveau
split \emph{stratifié-temporel} (270/60/45) a été assemblé dans
\texttt{data/datasets/ewat\_v4\_strat/} via \texttt{scripts/assemble\_dataset.py --stratified}.

\begin{table}[h]
\centering
\caption{Headline honnête Phase~B sur cible Chaos Mesh indépendante (ewat\_v4\_strat).}
\label{tab:archv2_headline}
\begin{tabular}{lcl}
\toprule
\textbf{Évaluation} & \textbf{macro-AUROC} & \textbf{IC 95\,\% bootstrap} \\
\midrule
B2 LR-OvR sur features brutes flatten (sans STGCN) & $\mathbf{0.920}$ & [0.878, 0.956] \\
B1 best (instance norm + last $k$, mean over $k$)  & $\mathbf{0.941}$ & [0.909, 0.970] \\
LOSO macro-AUROC (15 folds, retrain par fold)      & $0.930 \pm 0.007$ & --- \\
C1 STGCN end-to-end (instance norm + 15-way head)  & $0.863$ & [0.823, 0.905] \\
\bottomrule
\end{tabular}
\end{table}

\textbf{Le headline défensif final est donc $0.920$--$0.941$ sur cible Chaos Mesh}, IC explicite,
sans circularité. Notons que C1 (STGCN entraîné directement sur Chaos Mesh) reste sous la baseline
LR --- cohérent avec A5 : l'encodeur n'apporte pas de gain prédictif agrégé.
Le STGCN garde toutefois sa valeur géométrique (H1 silhouette $= 0.78$) et
ontologique (section~\ref{sec:ontology}).

\subsection{Phase C --- Précursion temporelle confirmée et open-set}

\paragraph{C2 --- A1 distant-window sur le modèle Chaos Mesh STGCN.}
Sur la version C1 du modèle (STGCN end-to-end + cible Chaos Mesh + v4\_strat) :

\begin{table}[h]
\centering
\caption{Distant-window sur STGCN+Chaos Mesh (v4\_strat, $k=6$, $n_{\text{test}}=45$).}
\label{tab:c2_distant}
\begin{tabular}{lcc}
\toprule
\textbf{Position fenêtre} & \textbf{macro-AUROC} & \textbf{IC 95\,\%} \\
\midrule
\texttt{last} (juste avant injection) & $\mathbf{0.876}$ & [0.838, 0.914] \\
\texttt{middle} (milieu régime normal) & $0.813$ & [0.763, 0.874] \\
\texttt{first} (début régime normal)  & $0.759$ & [0.708, 0.809] \\
\bottomrule
\end{tabular}
\end{table}

$\boldsymbol{\Delta($far\,$-$\,near$) = -0.116} \Rightarrow$ \textbf{précursion temporelle
réelle confirmée}. La dynamique pré-injection compte pour 12\,pp d'AUROC. La critique de fuite
soulevée par A1 sur labels EWAT s'inverse ici sur cible indépendante : EWAT exploite réellement
des patterns temporels qui s'amplifient à l'approche de l'injection.

\paragraph{C3 --- OpenMax open-set recognition.}
Pour répondre à la limitation A2 (LOSO top-1 $= 0$ sur scénario inédit avec OvR fermé), nous avons
implémenté OpenMax (Bendale \& Boult 2016, EVT/Weibull) dans \texttt{src/ewat/openset/openmax.py}
(11/11 tests unitaires). Évaluation LOSO complète (15 retrains $\times$ 60 époques) :
\texttt{experiments/architecture\_v2/openset\_eval.py}.

OpenMax permet d'attribuer une probabilité \emph{unknown} plutôt que de mal classifier
un scénario inédit dans une classe connue. C'est une réponse technique partielle à la critique
de généralisation, qui ouvre la voie à un système opérationnel capable de signaler "scénario
inconnu" plutôt que de produire une fausse alerte typée.

\paragraph{C5 --- H2 look-through retest sur ewat\_v4\_strat.}
Même avec des épisodes 2\,$\times$ plus longs (47--51 vs 21 steps), le mécanisme look-through
MMD$^2$ ne sépare toujours pas le drift bénin de l'anomalie ($p = 0.372$ paired t-test).
L'hypothèse "épisodes trop courts" n'explique pas tout : le mécanisme lui-même est
\textbf{fondamentalement falsifié}. C'est un résultat négatif robuste assumé.

\subsection{Synthèse}

\begin{itemize}
  \item Le chiffre $0.987$ (ou $0.973$ sur 5 graines) reporté en section~\ref{sec:h3}
    mesure la cohérence interne du clustering EWAT. Il est \textbf{circulaire} et ne doit pas
    être interprété comme métrique prédictive indépendante.
  \item Sur cible Chaos Mesh \textbf{indépendante} (ewat\_v4\_strat), le pipeline atteint
    $\mathbf{0.920}$ macro-AUROC [0.878, 0.956] avec un LR-OvR simple sur features brutes
    instance-normalisées, $0.941$ avec la meilleure configuration --- \textbf{c'est le headline
    défensif honnête}.
  \item La précursion temporelle est \textbf{réelle} ($\Delta($far\,$-$\,near$) = -0.116$ sur
    STGCN+Chaos Mesh) --- le modèle exploite la dynamique pré-injection, pas seulement la signature
    statique du scénario.
  \item L'open-set recognition (OpenMax/EVT) répond à la limitation de généralisation
    à un type inédit.
  \item H2 look-through est définitivement falsifié, indépendamment de la durée d'épisode.
\end{itemize}

\subsection{Pipeline opérationnel et budget de latence}
\label{sec:operational}

\paragraph{Pipeline opérationnel (Option B).}
Compte tenu de la neutralité agrégée du STGCN sur la cible indépendante
(section~\ref{sec:scenario_baseline}, A5 paired $\Delta(B_4-B_3)$ IC contient zéro
et C1 STGCN $= 0.863 <$ B2 LR $= 0.920$), le pipeline prédictif opérationnel
recommandé est :
\[
S(t) \rightarrow \text{instance norm} \rightarrow \text{LR-OvR 15-way} \rightarrow \text{softmax} \rightarrow \text{OpenMax/EVT}
\]
Le STGCN reste indispensable pour la \textbf{contribution géométrique} (H1
silhouette $= 0.782$ sur 10 graines ewat\_v3 avec config optimisée) et
\textbf{ontologique} (Phase 8, OWL/RDF formelle, 29 classes, HermiT cohérent).
La séparation entre la chaîne prédictive (LR sur features brutes) et la chaîne
de typage/ontologie (STGCN + siamois + clustering) clarifie le rôle de chaque
composant et évite la circularité d'évaluation H3.

\paragraph{Benchmark de latence (C-3, 200 itérations CPU).}
\begin{table}[h]
\centering
\caption{Latence end-to-end mesurée (ms). Budget formel : Étape 0 < 1 s,
Étape 1+2 < 2 s, Étape 3 < 1 s, Total < 5 s. Verdict global : 🟢 GREEN
(375$\times$ sous budget).}
\label{tab:latency}
\begin{tabular}{lcccr}
\toprule
\textbf{Étape} & \textbf{médiane} & \textbf{p95} & \textbf{p99} & \textbf{budget} \\
\midrule
Étape 0 (DriftDetector) & 0.01 & 0.01 & 0.02 & < 1000 ms \\
Étape 1+2 (encoder + siamois) & 0.96 & 1.97 & 2.76 & < 2000 ms \\
Étape 3 (précurseurs) & 1.85 & 3.91 & 4.74 & < 1000 ms \\
\textbf{TOTAL} & $\mathbf{9.26}$ & $\mathbf{13.28}$ & $\mathbf{17.64}$ & $\mathbf{< 5000}$ ms \\
\bottomrule
\end{tabular}
\end{table}

Le SLA latence opérationnelle est largement tenable sur CPU. À évaluer sur
GPU et sous charge concurrente (production réelle).

\paragraph{Power analysis a posteriori (C-5).}
Avec n\_test = 45 (ewat\_v3) et n\_test = 45-57 (ewat\_v4\_strat), seuls
\textbf{5/10 clusters} ont $n_{\text{pos}} \geq 5$ et sont statistiquement
reportables. La power moyenne sur ces 5 clusters est de 1.0 (Hanley-McNeil
SE) — les AUROCs reportés sont au-dessus du seuil $0.5$ avec puissance
parfaite. Pour atteindre une power $> 0.8$ sur un AUROC cible de $0.7$, il
faudrait $n_{\text{pos}} \geq 17$ par cluster — ewat\_v5+ devrait viser
$\geq 25$ répétitions par scénario.

\subsection{Limites résiduelles et travaux futurs}
\label{sec:residual_limits}

Au-delà des limites L1--L9 documentées en section~\ref{sec:limitations} et
en \texttt{docs/limitations.md}, l'audit méthodologique post-Phase~C identifie
huit limites résiduelles que nous reconnaissons et orientons en travaux futurs :

\begin{description}
  \item[L10 --- Surentraînement siamois sur ewat\_v4] H1 sil\_test = $0.467 \pm 0.156$
    sur 6 graines (vs $0.782 \pm 0.065$ sur v3), best\_epoch $\in [2, 7]$ vs $\sim$47.
    Cause probable : paires contrastives échantillonnées aléatoirement deviennent
    trop faciles sur le dataset plus diversifié.
    \textbf{Fix proposé} : hard-negative mining renforcé + curriculum learning.
  \item[L11 --- Latence] \textbf{Résolu} (Tableau~\ref{tab:latency}).
  \item[L12 --- OpenMax mitigé] Unknown AUROC $= 0.55 \pm 0.24$ (cible $0.7$).
    \textbf{Fix proposé} : alternatives Mahalanobis-OOD (Lee 2018), Energy-based
    (Liu 2020), ODIN (Liang 2017).
  \item[L13 --- Service graph $N=6$] Online Boutique vs production $100+$ services.
    Scalabilité non testée. \textbf{Fix} : benchmark à $N \geq 20$ services en travail futur.
  \item[L14 --- Validation cross-cluster] Seul \texttt{observit-cluster1}.
    \textbf{Fix} : Stratégie B fine-tuning + OTel Semantic Conventions strictes.
  \item[L15 --- Retraining cycle non défini] Pas de détection de drift de distribution
    déclenchant un re-entraînement. \textbf{Fix} : monitoring MLflow + détecteur d'OOD sur
    embeddings de production.
  \item[L16 --- 17 features hardcodées] Pas d'auto-feature-engineering. \textbf{Fix} :
    mRMR ou SHAP-based feature pruning.
  \item[L17 --- Ontologie sans audit SRE] Phase 8 validation formelle (HermiT cohérent,
    8/10 critères) mais pas d'audit expert opérationnel. \textbf{Fix} : 1 séance d'audit SRE
    + intégration runbook PagerDuty/OpsGenie.
\end{description}

Voir \texttt{docs/limitations.md} sections L10--L17 pour le détail diagnostique.

\subsection{Validation multi-seed finale (Phase~H + J + K, 2026-05-26)}
\label{sec:multiseed}

Suite à la critique de Phase~G (single seed $42$ donnant sil\_test $= 0.84$ et
$\Delta(\text{far}-\text{near}) = -0.05$), un sweep multi-seed sur 10 graines a
été mené sur \texttt{ewat\_v4\_strat}. \textbf{Verdict : Phase~G était un
outlier.}

\begin{table}[h]
\centering
\caption{Phase H — Multi-seed validation sur 10 graines (cible labels EWAT,
circulaire). $\Delta$(far--near) GENUINE\_PRECURSION sur 1/10 graines
uniquement — la graine 42 utilisée en Phase~G est un outlier.}
\label{tab:phase_h_multiseed}
\begin{tabular}{lcc}
\toprule
\textbf{Métrique} & \textbf{Mean $\pm$ Std} & \textbf{Range} \\
\midrule
H1 silhouette test         & $0.691 \pm 0.115$  & $[0.521, 0.839]$ \\
H3 AUROC peak (circulaire) & $0.990 \pm 0.012$  & $[0.959, 1.000]$ \\
A1 $\Delta$(far$-$near)    & $-0.012 \pm 0.022$ & $[-0.050, +0.019]$ \\
K\_optimal (silhouette)    & $11.8 \pm 2.1$     & $[9, 15]$ \\
\bottomrule
\end{tabular}
\end{table}

\begin{table}[h]
\centering
\caption{Phase J — Headline défensif sur cible Chaos Mesh indépendante.
Le LR-OvR est déterministe (solver lbfgs) → toutes les graines donnent
exactement le même chiffre. IC bootstrap est calculé séparément.}
\label{tab:phase_j_multiseed}
\begin{tabular}{lcc}
\toprule
\textbf{Métrique} & \textbf{Valeur déterministe} & \textbf{IC 95\,\% bootstrap} \\
\midrule
B2 stratified macro-AUROC & $\mathbf{0.9201}$ & $[0.878, 0.956]$ \\
B2 LOSO macro-AUROC       & $0.9298$          & --- \\
\bottomrule
\end{tabular}
\end{table}

\paragraph{Phase K --- diagnostics.}
\begin{itemize}
  \item \textbf{K.1 K-selection comparison} : silhouette vs Tibshirani gap
    statistic (Step~6 fix~6.4). Aucune méthode ne stabilise K — distribution
    [9, 15] vs [4, 12], agreement seulement 4/10 graines. K\_optimal est
    intrinsèquement instable sur n\_train $= 270$ (recommandation v5 :
    fixer $K = 10$ manuellement ou HDBSCAN).
  \item \textbf{K.3 Variance per-seed} : métriques stables (AUROC circulaire,
    B2 déterministe) vs métriques instables (sil\_test range $0.32$,
    K range $6$, A1 outlier seed~42). Diagramme :
    \texttt{experiments/multiseed/phase\_h/distribution.png}.
\end{itemize}

\paragraph{Verdict scientifique consolidé.}
\begin{enumerate}
  \item Le \textbf{headline défensif final} reste $\mathbf{B_2 = 0.9201}$
    [0.878, 0.956] sur Chaos Mesh \texttt{ewat\_v4\_strat} (déterministe,
    indépendant des labels EWAT).
  \item Le retrain Phase~G était un outlier (graine 42) --- Phase~H montre
    que $\text{sil} = 0.84$ et $\Delta = -0.05$ ne sont pas reproductibles.
  \item Les \textbf{38 fixes audit} (10-step plan) corrigent des bugs réels
    (NaN-aware \texttt{fit\_scaler}, \texttt{class\_weight='balanced'},
    instance norm exclusive, etc.) mais \emph{ne déplacent pas le headline
    défensif} ($B_2$ utilise un scaler local, pas l'encoder).
  \item Limites intrinsèques (L10, C-5, L13) : K\_optimal instable
    ($n_{\text{train}} = 270$ trop petit), surentraînement siamois
    (best\_epoch $\sim 3$ malgré semi-hard mining), $n_{\text{pos}} = 3$
    par scénario test (statistiquement faible).
\end{enumerate}

Détails complets dans \texttt{experiments/multiseed/phase\_h/} et
\texttt{experiments/multiseed/phase\_j/}.
